% !PRINTED 489
\weberpage{489}
Wenn wir die Indices der Variablen $u$, oder was damit gleichbedeutend ist, die $u$ selbst einer Permutation $\pi$ unterwerfen, so kann die Function $\psi$ sich ändern oder sie kann ungeändert bleiben. Bleibt sie ungeändert, so sagen wir, \gesperrt{sie gestatte die Permutation $\pi$}. Gestattet sie alle $\Pi(m)$ überhaupt möglichen Permutationen, so ist sie symmetrisch und kann durch die symmetrischen Grundfunctionen ausgedrückt werden (\S.\,44). Der andere extreme Fall ist der, dass die Function für alle $\Pi(m)$ Anordnungen lauter verschiedene Werthe annimmt. Im Allgemeinen werden gewisse Permutationen die Function $\psi$ ungeändert lassen, andere werden sie ändern. Wir stellen nun den Satz auf:

1. \gesperrt{Der Inbegriff aller der Permutationen der $u, u_1, u_2 \ldots u_{m-1}$, die eine ganze Function $\psi(u, u_1, u_2 \ldots u_{m-1})$ ungeändert lassen, ist eine Gruppe von Permutationen.}

Um ihn zu beweisen, nehmen wir an, es sei $\pi$ eine der Permutationen, die $\psi$ ungeändert lässt, was wir durch
\[
\psi = \psi_\pi
\]
ausdrücken. Da nun die $u$ unabhängige Variable sind, so muss diese identische Gleichung richtig bleiben, wenn die Variablen irgend einer Permutation $\pi'$ unterworfen werden. Wenn aber $\pi'$ auf $\psi_\pi$ angewandt wird, so ist das Ergebniss dasselbe, als wenn die zusammengesetzte Permutation $\pi\pi'$ auf $\psi$ angewandt wird. Es folgt also
\[
\psi_{\pi'} = \psi_{\pi\pi'}
\]
und wenn nun
\[
\psi = \psi_{\pi'}
\]
ist, so folgt
\[
\psi = \psi_{\pi\pi'},
\]
d.\,h. wenn die Function $\psi$ durch die Permutationen $\pi$ und $\pi'$ ungeändert bleibt, so bleibt sie auch durch $\pi\pi'$ ungeändert, wodurch nach der Definition der Gruppe der Satz, den wir ausgesprochen haben, bewiesen ist.

Wir haben darin ein Mittel, um Permutationsgruppen zu bilden, indem wir irgend eine Function von $m$ Variablen nehmen und alle Permutationen aufsuchen, die eine solche Function ungeändert lassen. Dass man auf diese Weise alle Permutationsgruppen bilden kann, wird sich nachher ergeben.
